% 2009 年北京卷文科第 16 题：四棱锥 P-ABCD，底面正方形，PD⊥底面，E 在棱 PB 上。
% 来源 PDF：普通高考/2009/2009北京文.pdf，SHA256 5b310eb0f16d1bd48553d20e7f16b10784e6c83787ef7dd8fa96e3065ef4c159
% 原页：第 1 页（A3 三栏）第 3 栏上部，印刷页 501。
% 重绘依据：底面正方形 A 左下、B 右下、C 右上、D 左上（斜投影），P 在 D 的正上方；
%   可见棱 PA、PB、PC、AB、BC 实线，被遮棱 PD、AD、DC 虚线；E 取 PB 中点，AE、CE 细线。
\begin{tikzpicture}[scale=0.95, semithick, every node/.style={inner sep=1.5pt}]
  \coordinate (A) at (0,0);
  \coordinate (B) at (2.4,0);
  \coordinate (C) at (3.2,1.1);
  \coordinate (D) at (0.8,1.1);
  \coordinate (P) at (0.8,3.1);
  % 被遮棱
  \draw[dashed] (P) -- (D);
  \draw[dashed] (A) -- (D);
  \draw[dashed] (D) -- (C);
  \draw[dashed] (A) -- (C);
  % 可见棱
  \draw (P) -- (A);
  \draw (P) -- (B);
  \draw (P) -- (C);
  \draw (A) -- (B) -- (C);
  % E 与 AE、CE
  \coordinate (E) at ($(P)!0.5!(B)$);
  \fill (E) circle (1.3pt);
  \draw (A) -- (E);
  \draw (C) -- (E);
  % 标注
  \node[anchor=north east] at (A) {$A$};
  \node[anchor=north west] at (B) {$B$};
  \node[anchor=west] at (C) {$C$};
  \node[anchor=east] at (D) {$D$};
  \node[anchor=south] at (P) {$P$};
  \node[anchor=west] at (E) {$E$};
\end{tikzpicture}
